% 底本：EGAIV.1.pdf，PDF第177--180页，原分册印刷第178--181页。
下面的命题是（21.7.3）的一个“形式”对应。

\egas{21.9.6}{命题}
设 $k_0\subseteq k\subseteq K_0$ 为三个特征 $p>0$ 的域，且 $[K_0:k]<\infty$。
置
\[
\begin{aligned}
 L_0&=K_0((T_1,\ldots,T_r)),&
 M_0&=K_0((T_1^p,\ldots,T_r^p)),\\
 L&=k((T_1,\ldots,T_r)),&
 M&=k((T_1^p,\ldots,T_r^p)),\\
 N&=k((T_1^{p^2},\ldots,T_r^{p^2})).
\end{aligned}
\]
设 $E/L_0$ 为有限型扩张。

（i）若 $\Upsilon_{K_0/k_0}$ 有限维，则
\begin{equation}
\begin{split}
 \rk_E\Omega^1_{E/M}-\rk_{K_0}\Omega^1_{K_0/k}
 =r+\trdeg_{L_0}E\qquad\qquad\\
 {}+d(E/K_0,k_0)+d(E/M_0,N/k_0).
\end{split}\tag{21.9.6.1}
\end{equation}
特别地，
\begin{equation}
 \rk_E\Omega^1_{E/M}-\rk_{K_0}\Omega^1_{K_0/k}
 \ge r+\trdeg_{L_0}E+d(E/K_0,k_0).
 \tag{21.9.6.2}
\end{equation}
若（21.9.6.2）取等号，则将 $k$ 换为任一满足 $k_0\subseteq k'\subseteq k$ 的域后，
该等号仍成立，只要所用域继续满足命题的有限性前提。
\cmt{L12-N06}{原文说可把 $k$ 换为任意 $k_0$ 与 $k$ 之间的子域。
此断言沿用命题开头 $[K_0:k]<\infty$ 的前提；若新 $k'$ 使 $[K_0:k']$ 无限，
（21.9.6.2）的两个秩未必有限，不能自动解释为一个定义良好的整数差。
正文把这一使用范围明确写出；下文单调性证明在有限秩范围内成立。}

（ii）再假设 $[k_0:k_0^p]<\infty$。
若 $(k_\alpha)$ 是 $K_0$ 的向下滤过的子域族，每个包含 $k_0$，
满足 $[K_0:k_\alpha]<\infty$，且
$\bigcap_\alpha k_\alpha(K_0^p)=k_0(K_0^p)$，
则存在 $\alpha$，使 $k=k_\alpha$ 时（21.9.6.2）取等号。

\eproof
（i）由（21.1.5，ii），在 $\Omega^1_{E/M}$ 中把 $M$ 换成 $M(L_0^p)$ 不改变这个模。
由于 $L_0^p=K_0^p((T_1^p,\ldots,T_r^p))$，且 $[k(K_0^p):k]<\infty$，
\[
 M(L_0^p)=k(K_0^p)((T_1^p,\ldots,T_r^p)).
\]
同样，把 $k$ 换成 $k(K_0^p)$ 不改变 $\Omega^1_{K_0/k}$。
所以可假设 $K_0^p\subseteq k$，以下一直采用这一假设。
\cmt{L12-G13}{这一替换还需要检查（21.9.6.1）右端的缺陷项。
替换后 $N$ 成为 $N(M_0^p)$：利用 $[k(K_0^p):k]<\infty$ 可逐系数展开得到此等式。
由（21.1.5.1），$\Upsilon_{E/M_0/N}$ 不变，故
$d(E/M_0,N/k_0)=\dim_E\Ker(\Upsilon_{E/M_0/k_0}\to\Upsilon_{E/M_0/N})$ 不变。
正文归约因而同时保持等式两端；这补明原文略去的一个相容性检查。}
再引入
\[
 k'_0=k_0(K_0^p),\qquad N'=k'_0((T_1^{p^2},\ldots,T_r^{p^2})),
\]
于是有域的交换图
\[
\begin{tikzcd}[column sep=small,row sep=large]
 &K_0\arrow[rr]&&M_0\arrow[r]&L_0\arrow[r]&E\\
 &k\arrow[u]\arrow[r]&N\arrow[r]&M\arrow[u]\arrow[r]&L\arrow[u]&\\
 k_0\arrow[r]&k'_0\arrow[u]\arrow[r]&N'\arrow[u]&&&
\end{tikzcd}
\]
以 $\delta$ 表示（21.9.6.2）左端减去右端所得的差。
由（21.7.3），
\[
 \rk_E\Omega^1_{E/M}-\rk_{M_0}\Omega^1_{M_0/M}
 =\trdeg_{M_0}E+d(E/M_0,M).
\]
\cmt{L12-T07}{底本第179页在首次应用（21.7.3）的这一行把第二个微分模的分子、分母写反，
成为 $\Omega^1_{M/M_0}$；但实际包含关系是 $M\subseteq M_0$。
正文改为 $\Omega^1_{M_0/M}$，这也与原文紧接着的下一条秩等式一致。}
$L_0/M_0$ 有限，故 $\trdeg_{M_0}E=\trdeg_{L_0}E$。
又因 $K_0$ 在 $k=k(K_0^p)$ 上的一组 $p$-基也是 $M_0/M$ 的一组 $p$-基，
由（21.4.5），
\[
 \rk_{M_0}\Omega^1_{M_0/M}=\rk_{K_0}\Omega^1_{K_0/k}.
\]
于是
\begin{equation}
 \delta=d(E/M_0,M)-d(E/K_0,k_0)-r.
 \tag{21.9.6.3}
\end{equation}
先回忆下面的经典引理。

\egas{21.9.6.4}{引理}
对任意域 $k$，形式幂级数环的分式域
$K=k((T_1,\ldots,T_r))$ 是 $k$ 的可分扩张。

\eproof
为使论述完整，简要回忆证明。
只须对每个有限扩张 $k'/k$ 证明 $K\otimes_kk'$ 无非零幂零元
（Bourbaki，\textit{Alg.}，第VIII章，§7，第3节，定理1的证明）。
置 $A=k[[T_1,\ldots,T_r]]$、$S=A\setminus\{0\}$，则
\[
 K\otimes_kk'=S^{-1}(A\otimes_kk').
\]
因为 $k'/k$ 有限，$A\otimes_kk'$ 典范同构于整环
$A'=k'[[T_1,\ldots,T_r]]$，且 $A$ 识别为 $A'$ 的子环。
相应局部化仍为整环，因而约化。
\finishproof

\noindent\textit{命题（21.9.6）的证明续。}
由此引理和（21.6.2，ii），$d(L_0/K_0,k_0)=0$，
再由（21.6.4.4）得到
\[
 d(E/K_0,k_0)=d(E/L_0,k_0).
\]
另一方面，（21.6.4.4）、（21.6.4.2）分别给出
\begin{align*}
 d(E/M_0,M)&=d(E/L_0,M)+d(L_0/M_0,M),\\
 d(L_0/M_0,M)&=d(L_0/M_0,M/N)+d(L_0/M_0,N).
\end{align*}
所以
\begin{align*}
 \delta={}&d(E/L_0,M)+d(L_0/M_0,N)-d(E/L_0,k_0)\\
          &+d(L_0/M_0,M/N)-r.
\end{align*}
下面证明
\begin{equation}
 d(L_0/M_0,M/N)=r.
 \tag{21.9.6.5}
\end{equation}
按定义，
\[
 d(L_0/M_0,M/N)
 =\rk_{L_0}\Upsilon_{L_0/M/N}-\rk_{M_0}\Upsilon_{M_0/M/N}.
\]
考虑两条正合列
\[
\begin{gathered}
 0\longrightarrow\Upsilon_{L_0/M/N}
 \longrightarrow\Omega^1_{M/N}\otimes_ML_0
 \longrightarrow\Omega^1_{L_0/N}
 \longrightarrow\Omega^1_{L_0/M}\longrightarrow0,\\
 0\longrightarrow\Upsilon_{M_0/M/N}
 \longrightarrow\Omega^1_{M/N}\otimes_MM_0
 \longrightarrow\Omega^1_{M_0/N}
 \longrightarrow\Omega^1_{M_0/M}\longrightarrow0.
\end{gathered}
\]
交错秩相减，得到
\begin{align*}
 d(L_0/M_0,M/N)
 ={}&\rk_{L_0}\Omega^1_{L_0/M}-\rk_{L_0}\Omega^1_{L_0/N}\\
 &-\rk_{M_0}\Omega^1_{M_0/M}+\rk_{M_0}\Omega^1_{M_0/N}.
\end{align*}
而 $N(L_0^p)\supseteq M$，故
$\Omega^1_{L_0/N}=\Omega^1_{L_0/M}$（21.1.5，ii）。
此外 $T_1^p,\ldots,T_r^p$ 构成 $M/N$ 的 $p$-基，
$K_0/k$ 的一组 $p$-基又是 $M_0/M$ 的一组 $p$-基；
因 $M_0^p\subseteq N$，由（21.4.5）和（21.1.10）得
\[
 \rk_{M_0}\Omega^1_{M_0/N}
 =r+\rk_{K_0}\Omega^1_{K_0/k}
 =r+\rk_{M_0}\Omega^1_{M_0/M}.
\]
这就证明了（21.9.6.5）。

再由（21.6.4.2），
\begin{align*}
 d(E/L_0,M)-d(E/L_0,k_0)&=d(E/L_0,M/k_0),\\
 d(L_0/M_0,N)&=d(L_0/M_0,N/k)+d(L_0/M_0,k).
\end{align*}
$L_0/K_0$ 可分（21.9.6.4），所以 $d(L_0/K_0,k)=0$，
进而 $d(L_0/M_0,k)=0$（21.6.2及21.6.4）。因此
\[
 \delta=d(E/L_0,M/k_0)+d(L_0/M_0,N/k).
\]
对第一项连续应用两次（21.6.4.2），得
\begin{align*}
 d(E/L_0,M/k_0)
 ={}&d(E/L_0,M/N)+d(E/L_0,N/k)\\
 &+d(E/L_0,k/k_0).
\end{align*}
因 $M\subseteq N(L_0^p)$，有 $d(E/L_0,M/N)=0$（21.6.2）；
再用（21.6.4.4）合并另外两项，得到
\[
 \delta=d(E/L_0,k/k_0)+d(E/M_0,N/k).
\]
$M_0$、$L_0$ 都在 $K_0$ 上可分（21.9.6.4），
故 $d(L_0/K_0,k/k_0)=d(M_0/K_0,k/k_0)=0$（21.6.2）。
再应用（21.6.4.4），
\[
 d(E/L_0,k/k_0)=d(E/K_0,k/k_0)=d(E/M_0,k/k_0).
\]
最后用（21.6.4.2），即得
\[
 \delta=d(E/M_0,N/k_0)\ge0.
\]
这证明（21.9.6.1）及（21.9.6.2）。
若将 $k$ 换成子域 $k'$，则 $N$ 换成子域 $N''$，而
$d(E/M_0,N''/k_0)\le d(E/M_0,N/k_0)$（21.6.4.2）。
所以在秩差仍有定义的上述范围内，取等号的性质保留，证明了（i）的最后一项。

（ii）仍可先把 $k_\alpha$ 换成 $k_\alpha(K_0^p)$，故假设 $K_0^p\subseteq k_\alpha$。
置 $N_\alpha=k_\alpha((T_1^{p^2},\ldots,T_r^{p^2}))$。
由（21.9.6.1）和（21.8.3），只须证明
\[
 \bigcap_\alpha N_\alpha=k_0(M_0^p),
\]
因为 $E/M_0$ 有限型，且每个 $N_\alpha$ 均包含 $M_0^p$。
由（21.8.8.3）及假设 $\bigcap_\alpha k_\alpha=k_0(K_0^p)$，有
\begin{equation}
 \bigcap_\alpha N_\alpha
 =N'=(k_0(K_0^p))((T_1^{p^2},\ldots,T_r^{p^2})).
 \tag{21.9.6.6}
\end{equation}
又 $k_0(M_0^p)=k_0(K_0^p((T_1^{p^2},\ldots,T_r^{p^2})))$。
所以要完成（ii）的证明，只须使用下一个引理。

\egas{21.9.6.7}{引理}
设 $k_0$ 为特征 $p>0$ 的域，且 $[k_0:k_0^p]<\infty$；$K_0/k_0$ 为域扩张。
则
\begin{equation}
 (k_0(K_0^p))((T_1,\ldots,T_r))
 =k_0(K_0^p((T_1,\ldots,T_r))).
 \tag{21.9.6.8}
\end{equation}

\eproof
由于 $k_0^p\subseteq K_0^p$，
\[
 k_0^p(K_0^p((T_1,\ldots,T_r)))=K_0^p((T_1,\ldots,T_r)).
\]
若 $c_1,\ldots,c_n$ 是 $k_0/k_0^p$ 的一组基，
则它们同时生成（21.9.6.8）两端作为
$K_0^p((T_1,\ldots,T_r))$-向量空间，故两端相等。
\cmt{L12-N07}{这里用的是有限生成系，不要求 $(c_i)$ 在扩大系数域后仍线性无关。
可先逐系数将左端的幂级数写成有限和 $\sum_i c_i f_i$，再取分式域。
有限性 $[k_0:k_0^p]<\infty$ 使同一有限组 $(c_i)$ 对所有系数同时适用；
下一条原著反例正说明不能直接删除这一假设。}
\finishproof

由此引理，（21.9.6.6）右端就是 $k_0(M_0^p)$，
因而（21.8.3）保证某个 $N_\alpha$ 对 $E/M_0$ 为 $k_0$-可容许域，
即相应的 $\delta=0$；命题（21.9.6）证毕。
